% NOTE: In this file all the ToC entries are manually added because they
% ```   dont show up on the ToCs with the * method otherwise
\thispagestyle{empty}
\begin{center}
    \Huge{
        \textbf{
            \color{red}
            JavaHotel
            \includegraphics[scale=2]{java-hotel}
            \color{black}
            \vspace{5mm}
        }
    }
    \\
    \large{
        Ένα πρόγραμμα για την υλοποίηση ενός συστήματος κρατήσεων, ακυρώσεων
        και ειδοποιήσεων σε ένα ξενοδοχείο.
    }
    \vspace{5mm}\\
  %  \includesvg[scale=5]{../svg/cover}\vspace{5mm}\\
    \small{\emph{
    Διονύσης Νικολόπουλος, ice18390126@uniwa.gr
    \\
    Τμήμα Μηχανικών Πληροφορικής και Υπολογιστών
    \\
    Πανεπιστήμιο Δυτικής Αττικής, Ελλάδα, Αθήνα
    \\
    Μάιος - Ιούνιος 2022
    \\
    }}
    \vspace*{5mm}
\end{center}
\footnotesize
Keywords:
\emph{Γλώσσα προγραμματισμού Java, Remote Method Invocation, Κατανεμημένα
Συστήματα}
\normalsize


\section*{Θέμα και Λεπτομέρειες Υλοποίησης}
\begin{addmargin}[2em]{3em}% 2em left, 3em right
\footnotesize{
    Ζητούμενο της εργασίας ήταν η κατασκευή ενός προγράμματος τέλεσης 5 πιθανών
    διαδικασιών:
    \begin{itemize}
        \item Την κράτηση ενός αριθμού και τύπου δωματίων στο όνομα ενός πελάτη.
        \item Την ακύρωση ενός αριθμού προηγουμένως κρατημένων δωματιών.
        \item Την εμφάνιση των διαθέσιμων δωματιών του ξενεδοχείου.
        \item Την εμφάνιση της λίστας κρατήσεων του ξενεδοχείου.
        \item Την ειδοποίηση πελατών σε περίπτωση που ένας αριθμός του
            επιλεγμένου τύπου δωματίων γίνει διαθέσιμος.
    \end{itemize}
    Όλα τα ερωτήματα της εκφώνησης, όπως αυτή φαίνεται στον φάκελο exc/ του
    repository της εργασίας, έχουν πραγματοποιηθεί.
    \\
    Το πρόγραμμα δημιουργήθηκε και δοκιμάστηκε σε περιβάλλον Arch Linux
    με kernel 5.18.1-arch1-1, με OpenJDK version 11.0.15.
    \\
    Ο κώδικας που γράφθηκε υπόκειται στην άδεια GNU GPLv3.
    \\
    Η Εργασία αυτή πραγματοποιήθηκε στο πλαίσιο του μαθήματος των Κατανεμημένων
    Συστημάτων του προγράμματος σπουδών Πανεπιστημίου Δυτικής Αττικής, τμήματος
    Μηχανικών και Πληροφορικής Υπολογιστών.
    \\
    \color{blue}
    Κάθε αρχείο αυτής της εργασίας, συμπεριλαμβανομένου και του παρόντος
    εγχειριδίου βρίσκονται αποθηκευμένα σε ένα αποθετήριο στο\gitlab{}\ \  \ 
    GitLab,
    \href{https://gitlab.com/nionios/java-hotel}
    {στον υπερσύνδεσμο αυτό.}
    \color{black}
    Εκεί φαίνεται αναλυτικά το πως και για πόσο έγινε η δουλειά πάνω
    στο πρόγραμμα.
}
\end{addmargin}
